\scanpage{46}
% archive foliation: 45

\deleted{sont \uline{\unsure{canoniquement}} équivalentes, \unsure{Si on les} \unsure{couples} essentiellement identifiés entre eux les différents $(C,\alpha)$, on trouve ainsi une catégorie de Deuring \uline{canonique}}
\inserted{l'équivalence canonique à isom près}
sont équivalents: la catégorie des équivalences de $(C,\alpha)$ vers $(C',\alpha')$ est un groupoïde connexe de $\pi_1\quad\mathbb Z/2\mathbb Z$:
i.e. tt d'abord les équivalences sont isom., et si on a deux équiv. $\varphi,\varphi'$, il y a \uline{au signe près} un seul isom. entre $\varphi$ et $\varphi'$.

Si on se donne maintenant un couple $(C,\alpha)$, il lui est associé une \deleted{torseur sur} gerbe sur
$\mathbb A=\prod_\ell\mathbb Z_\ell$ \inserted{(y compris $\ell=p$, $\ell=\infty$)}
\inserted{de groupe $\mu_2$,}
savoir celle des équivalences \inserted{de $(C,T)$} avec le \deleted{gerbe} couple type de $(C(\mathbb A),T)$, qui en la place $p$ soit compatible avec la donnée de $\alpha$.
Cette gerbe, $\unsure{\mathfrak g}_{(C,\alpha)}$ est déterminée, ``à un signe près'' \unsure{à} équivalence unique près, indépendamment du choix de $(C,\alpha)$ (puisqu'il en est ainsi de $(C,\alpha)$ lui-même), \deleted{et}
\deleted{$\unsure{\mathcal T}=\unsure{\mathcal T}(\overline{\mathbb F}_p/\mathbb F_p)$ \unsure{opère trivialement}.}
mais elle dépend à priori \inserted{\unsure{$\simeq$ pour $(C,\alpha)$ fixé},} du choix d'un $\overline{\mathbb F}_p$
\deleted{(ou du moins de $\mathbb F_{p^2}$)}
--- tout comme la structure $\alpha$ elle-même!

\margin{\deleted{\unsure{Cristal} \unsure{isocline} de rang $2$, pente $\frac12$ sur $\overline{\mathbb F}_p$, en descendant \uline{canoniquement} à $\mathbb F_{p^2}$.}}

On peut s'en tirer en définissant \uline{canoniquement} un $\mathbb F_{p^2}$ (ou même un $\overline{\mathbb F}_p$, si on y tenait \ldots) par les polynômes cyclotomiques $\Phi_{p^n-1}$,
\[
\overline{\mathbb F}_p=
\underset{n}{\underrightarrow{\lim}}\,
\bigl(\mathbb F_{p^n}=\mathbb F_p[T]/\Phi_{p^n-1}(T)\bigr)
\qquad\ldots
\]
